\documentclass[sigplan,10pt,anonymous,review]{acmart}
\settopmatter{printfolios=true,printccs=false,printacmref=false}
\acmConference[CPP '27]{Certified Programs and Proofs}{January 2027}{}
\acmYear{2027}\copyrightyear{2027}
\settopmatter{printacmref=false}
\renewcommand\footnotetextcopyrightpermission[1]{}
\pagestyle{empty}

\begin{document}
\title{Tower Induction and Interaction Trees: Experience Report}

\begin{abstract}
Parameterized coinduction, or \texttt{paco}, improves on Rocq's awkward
\texttt{cofix} tactic for low-level coinductive reasoning. But \texttt{paco}
still relies on many interlocking definitions and a wide array of library
tactics that can make proofs complex. The \texttt{rocq-coinduction} library,
which mechanizes the tower induction principle of Sch\"afer and Smolka, promises
shorter and more uniform proofs: tower induction is the single proof method for
both coinduction and soundness for up-to techniques, and it requires fewer
tactics than \texttt{paco} to employ.

We assessed that promise at scale by porting the coinductive-proof
infrastructure of Interaction Trees (ITrees), a large coinduction-heavy library,
from \texttt{paco} to \texttt{rocq-coinduction}. We found a mixed result: proof
text decreased by 12\%, but tactic text nearly tripled.
\texttt{rocq-coinduction} required significant client-side plumbing around its
exposed tactics and contained several usability bugs that required verbose
workarounds. Development also stalled because \texttt{rocq-coinduction}'s
tactics were written in OCaml over complex Rocq infrastructure, which made them
difficult to debug.

Observing that \texttt{rocq-coinduction} was theoretically rich but awkward in
implementation, we rewrote its tactic layer from scratch. The new implementation
is Rocq-native, replacing the OCaml plugin with Gallina and Ltac, and is 65\%
smaller than the code it supersedes while exporting more expressive tactics and
removing the limitations above. Rerunning the ITrees port against the
reimplemented library keeps the proof reduction and takes tactic text from
178\% above the \texttt{paco} baseline to 44.5\% below it. These results suggest
that the re-engineered tower-based coinduction library is not only a theoretical
simplification but a practical improvement over \texttt{paco} for large
coinductive developments.
\end{abstract}

\maketitle
\thispagestyle{empty}
\end{document}
